\documentclass{article} 

\usepackage[utf8]{inputenc} 
\usepackage[T1]{fontenc}    
\usepackage[english]{babel} 
\usepackage{amsmath}        
\usepackage{amssymb}        
\usepackage{amsfonts}       
\usepackage{graphicx}       
\usepackage{hyperref}       
\usepackage{geometry}       
\usepackage{enumitem}       
\usepackage{titling}
\usepackage{braket}
\usepackage{graphicx}
\usepackage{multicol}
\usepackage{lipsum}

% Page layout
\geometry{margin=1cm}      
\geometry{bottom=2cm}
\geometry{top=0.5cm}
\setlength{\parindent}{0pt} 
\setlength{\parskip}{6pt}   % Space between paragraphs
% math commands
\newcommand{\Lagrange}{\mathcal{L}}
\newcommand{\Hamilton}{\mathcal{H}}
\newcommand{\deriv}[3][]{\frac{d^{#1} #2}{d #3^{#1}}}
\newcommand{\pderiv}[3][]{\frac{\partial^{#1} #2}{\partial #3^{#1}}}

\renewcommand{\labelenumii}{\theenumii}
\renewcommand{\theenumii}{\theenumi.\arabic{enumii}.}

\graphicspath{{../plots}}

\pretitle{\begin{center}\LARGE}
\posttitle{\par\end{center}\vskip 0.5em}

\title{NDT Homework calculations}
\author{Capt. A.B.}
\date{\today}

\begin{document}
\maketitle

\section{Math used}
\subsection{Moment of radial redistribition}
Let \(\chi(r)\) be the radial redistribution function
\[
  \chi(r) = 1 + D\left[\ e^{-2\alpha(r - r^*)/R}  - 2e^{-\alpha(r-r^*)/R} \right]
\]
The integral function of the first radial moment is
\begin{align*}
  I(r) &= \int r\chi(r)dr \\
       &= \frac{r^2}{2}  - \frac{RD}{2\alpha}re^{-2\alpha(r - r^*)/R} - \frac{R^2D}{4\alpha^2}e^{-2\alpha(r - r^*)/R} + 2\frac{RD}{\alpha}re^{-\alpha(r - r^*)/R} + 2\frac{R^2D}{\alpha^2}e^{-\alpha(r - r^*)/R} + C
\end{align*}
Therefore for the fundamental theorem 
\[
  \int_{r_0}^R r\chi(r)dr = I(R) - I(r_0)
\]
\subsection{Duhamel method}
Let $u(r,t)$ be the solution of the diffusion problem
$$\[
  \begin{cases}
    \pderiv{u}{t} = D \nabla^2 u + f(r,t)  \quad & r \in (0, R), t > 0 \\
    u(r,0) = 0 \quad & r \in [0, R], t = 0 \\
    + BC
  \end{cases}
\]

and \(w(r,t,s)\) family of solutions with \(s \in [0,t]\) of the homogeneous problem
\[
  \begin{cases}
    \pderiv{w}{t} = D \nabla^2 w  \quad & r \in (0, R), t > 0 \\
    w(r,0,s) = f(r,s) \quad & r \in [0, R], t = 0 \\
    + BC
  \end{cases} 
\]

the solution of the non-homogeneous problem is given by the Duhamel formula
\[
  u(r,t) = \int_0^t w(r,t-s,s) \, ds
\]
\section{Thermal analysis up to fuel outer radius}
We consider fixed channel geometry. Coolant passage area [m$^2$]:
\[
  A_p = \frac{\sqrt{3}}{2}p^2 - 0.5\pi R_{co}^2
\]
where $p$ pitch, $R_{co}$ cladding outer radius.
Coolant axial temperature profile by solving (f = 1/2 channel factor)
\[ 
  \pderiv{T_c}{z} = f\frac{q'(z)}{\rhoG c_p(T_c) }
\]
Convective HTC [W/m$^2$K]
\begin{align*}
  & v = v(z) = \frac{G}{\rho A_p} \\
  & Nu(z) = 7 + 0.025Pe(z)^{0.8} \qquad Pe = \frac{\rho v D_h c_p}{k}\\
  & h(z) = \frac{k}{D_h}Nu(z)
\end{align*}

Cladding inner and outer temperature axial profiles with electrical analogy
\begin{align*}
  & T_{cl,o}(z) = T_{cool}(z) + f\frac{q'(z)}{2\pi R_{co}h(z)}\\
  & T_{cl,i}(z) = T_{cl,o}(z) + \frac{q'(z)}{2\pi k_{cl}H_a}\ln\left(\frac{R_{co}}{R_{ci}}\right)
\end{align*}
where $k_{cl} = k_{cl}[0.5(T_{cl,i} + T_{cl,o})]$. To get the fuel outer temperature along z, we apply again the electrical analogy across the gap
\[
  T_{fo}(z) = T_{cl,i}(z) + q'(z)\frac{t_{gap}}{2\pi k_{gap}H_a}
\]
where $k_{gap} = k_{gap}[0.5(T_{cl,i} + T_{fo})]$. 

\section{Thermal analysis inside the fuel}
\subsection{Conductivity integral}
The fuel thermal conductivity is given by the correlation 
\begin{align*}
  & k = k(T, y_{Pu}, x, p, \beta) =  1.755 + (k_0 - 1.755)e^{-\beta/128.75} \quad [W m^{-1} K^{-1}]\\
  & k_0 = \left[\frac{1}{A + BT} + \frac{D}{T^2}e^{-E/T} \right](1 - p)^{2.5}
\end{align*}
Integrating in temeprature range \([T_o, T_i]\), which corresponds to radial range \([r_i,r_o]\) we obtain
\[
  \int_{T_o}^{T_i} k dT = 1.755(T_i -  T_o)(1 - e^{-\beta/128.75}) + \left[ \frac{1}{B}\log\left(\frac{A + BT_i}{A + BT_o}\right) + \frac{D}{E}(e^{-E/T_i} - e^{-E/T_o})\right](1-p)^{2.5}e^{-\beta/128.75}
\]
The values of the properties \(y_{Pu}, x, p, \beta\), if not constant, can be averaged in \([r_i,r_o]\)

\subsection{Plutonium redistribution}
The Pu radial redistribution is given by the empirical formula
\begin{equation}
  \chi(r) = \frac{y(r)}{y_0} =  1 + D\left[\ e^{-2\alpha(r - r^*)/R}  - 2e^{-\alpha(r-r^*)/R} \right] 
  \label{eq:radial_red}
\end{equation}
\[
\]
where $D$ = 0.01, $\alpha = 10$ and  $r^*$ is chosen to conserve the Pu mass. Considering a pellet with central void of radius \(R_0\) the initial uniform fraction is given by
\[
  y_0 = 0.29
\]
The mass balance imposes
\[
  \int_V y(r) dV = y_0V \quad \Rightarrow \quad 2\pi H\int_{r_0}^R r \chi(r)dr = V
\]
Let \(\xi = r/R\), \(d\xi = Rdr\), therefore
\[
  V + 2\pi R^2HD\;e^{2\alpha\xi^*}\int_{\xi_0}^1 \xi e^{-2\alpha\xi}d\xi - 4\pi R^2HD\;e^{\alpha\xi^*}\int_{\xi_0}^1 \xi e^{-\alpha\xi}d\xi = V
\]
\[
  e^{2\alpha\xi^*}\int_{\xi_0}^1 \xi e^{-2\alpha\xi}d\xi - 2e^{\alpha\xi^*}\int_{\xi_0}^1 \xi e^{-\alpha\xi}d\xi = 0
\]
Defining the integral
\[
  J_n = \int_{\xi_0}^1 \xi e^{-n\xi}d\xi = \frac{1}{n\alpha}\left( \xi_0e^{-n\alpha\xi_0} - e^{-n\alpha} \right) + \frac{1}{n^2\alpha^2}\left(e^{-n\alpha\xi_0} - e^{-n\alpha} \right)
\]

We obtain an algebraic equation for \(\xi^*\)
\begin{align*}
  & J_2e^{2\alpha\xi^*} - 2J_1e^{\alpha\xi^*} = 0\\
  & e^{\alpha\xi^*} = \frac{2J_1}{J_2}  \quad \Rightarrow \quad \xi^* = \frac{1}{\alpha}\ln\left(\frac{2J_1}{J_2}\right) \quad \Rightarrow \quad r^* = \frac{R}{\alpha}\ln\left(\frac{2J_1}{J_2}\right)
\end{align*}

The power density is considered to be proportional to the Pu fraction
\[
  q'''(r,z) = q'''_0\, \chi  = \frac{ q'_0(z) }{\Omega}\chi(r)
\]

\subsection{Fuel temperature profile}
The fuel temperature profile is obtained by solving numerically the heat conduction equation at each axial node
\[
   \frac{1}{r}\pderiv{}{r}\left(r k \pderiv{T}{r}\right) = -q'''\\
\]
\[
  rk\pderiv{T}{r} = -q'''_0(z)\int r\chi(r)dr = -q'''_0(z)I(r) + C(z)
\]
The boundary conditions are
\[
  \begin{cases}
    T(R) = T_{fo} \\
    \pderiv{T}{r}(r_0) = 0
  \end{cases}
\]

Applying the no flux condition at the center we obtain
\[
  C(z) = q'''_0(z)I(r_0)
\]
We now have an initual value problem which can be solved numerically at each axial node
\begin{equation}
  \pderiv{T}{r} = -q'''_0(z)\frac{I(r)}{r\,k(T)} + \frac{C(z)}{r\,k(T)}
  \label{eq:IVP_fuel_T}
\end{equation}

\subsection{Restructured fuel}
In restructured fuel we can subdivide in 3 regions
\begin{enumerate}
  \item Columnar grain region: \(r_0 < r < r_1\);  \(T_0 < T < T_1 = 1800\) °C;  \(\rho_1 = 0.95 - 0.99\%\; TD\)\\
  \item Equiaxed grain region: \(r_1 < r < r_2\);  \(T_1 < T < T_2 = 1600\) °C;  \(\rho_2 = 0.95 - 0.97\%\; TD\)\\
  \item As-fabricated region: \(r_2 < r < R\);  \(T_2 < T < T_{fo}\);  \(\rho_3 = \rho_s\)\\
\end{enumerate}

Porosity of each region is given by
\[
  p_i = 1 - \frac{\rho_i}{TD}
\]
% In region 1 and 2 the linear power densities change due to densification and Pu redistribition
% \[
%   q'_i = q'_3\frac{\rho_i}{\rho_3}\frac{\bar{y}_i}{y_{i0}} \quad i = 1,2 \quad \text{(2 probably Pu negligible)}
% \]

The unknowns of the problem are the radii \(r_0, r_1, r_2\) and the temperature \(T_0\). By choosing a guess radius \(r_0\) we can compute the Pu and power redistribition. The power density is given by
\[
  q'''(r,z) = \frac{q'''_0(z)\chi(r)}{\Omega} = \frac{q'_0(z)\chi(r)}{\pi(R^2 - r_o^2)}
\]

We can solve the IVP defined by equation \ref{eq:IVP_fuel_T}, up until \(r_0\). By doing so owe find \(T_0\) and \(r_1\), \(r_2\), which are at the radial nodes where the temperature is respectively \(T_1\) and \(T_2\). A new value of \(r_0\) is given by the mass balance
\begin{equation} 
  r_0^2 = \left(\frac{\rho_1 - \rho_2}{\rho_1}\right)r_1^2 - \left(\frac{\rho_2 - \rho_3}{\rho_1}\right)r_2^2
  \label{eq:mass_balance}
\end{equation}

Repeat until convergence.



% While integrating the conduction equation in the three regions gives
% \begin{align}
%   & \int_{T_1}^{T_0}kdT = \frac{q'_1}{4\pi R^2}(r_1^2 - r_0^2) - C_1\log\left(\frac{r_1}{r_0}\right)\\ 
%   & \int_{T_2}^{T_1}kdT = \frac{q'_2}{4\pi R^2}(r_2^2 - r_1^2) - C_2\log\left(\frac{r_2}{r_1}\right)\\ 
%   & \int_{T_1}^{T_0}kdT = \frac{q'_1}{4\pi}\left(1 - \frac{r_2^2}{R^2}\right) - C_3\log\left(\frac{r_1}{r_0}\right)
% \end{align}
% where the constants \(C_i\) are evaluated by the boundary conditions.
% \[
%   \begin{cases}
%     & k_3\deriv{T}{r}(r_2) = k_2\deriv{T}{r}{r_2}\\
%     & k_2\deriv{T}{r}(r_1) = k_1\deriv{T}{r}{r_1}\\
%     & \deriv{T}{r}(r_0) = 0
%   \end{cases}
%
% Applying the boundary conditions we obtain
% \[
%   \begin{cases}
%   & C_1 = \frac{q'_1}{2\pi R^2}r_0^2\\
%   & C_2 = \frac{1}{2}\left[ q'_2r_1^2 + q'_1(r_0^2 - r_1^2) \right]\\
%   & C_3 = 0
%   \end{cases}
% \]
% Computing the conductivity integral as shown in section 2.1 equations (1) to (4) become a system of non-linear equations in \(r_0, r_1, r_2\) and \(T_0\), which can be solved numerically.


\section{Fission gas release}

\subsection{Production}
The production rate of gas atoms with fission yield $y$ is given by
\[
  \deriv{p}{t} = y\dot{F}V_g
\]
where $\dot{F}$ is the fission rate and $V_g$ is the volume of the grain. Integrating we obtain
\[
  p(t) = y\dot{F}V_g\:t \quad \text{[atoms]}
\]

\subsection{Retention in the grain}

The diffusion problem for a spherical grain is
\[
\begin{cases}
  \pderiv{u}{t} = D \frac{1}{r^2}\pderiv{}{r}( r^2 \pderiv{u}{r} ) + y\dot{F}  \quad & r \in (0, R), t > 0 \\
  u(r,0) = 0 \quad & r \in [0, R], t = 0 \\
  \pderiv{u}{r}(0,t) = 0 \quad & t > 0\\
  u(R, t) = 0 \quad & t > 0
\end{cases}
\]

where $u(r,t)$ is the concentration of gas atoms in the grain. The homogeneous problem is 
\[
\begin{cases}
  \pderiv{w}{t} = D \frac{1}{r^2}\pderiv{}{r}( r^2 \pderiv{w}{r} )  \quad & r \in (0, R), t > 0 \\
  w(r,0) = y\dot{F} \quad & r \in [0, R], t = 0 \\
  \pderiv{w}{r}(0,t) = 0 \quad & t > 0\\
  w(R, t) = 0 \quad & t > 0
  \end{cases}
\]

apply the separation of variables $w(r,t) = \theta(t)\phi(r)$
\[
  \frac{1}{D\theta}\pderiv{\theta}{t} = \frac{1}{r^2\phi}\pderiv{}{r}( r^2 \pderiv{\phi}{r} ) = -k^2
\]

the time part is
\[
  \theta(t) = Ce^{-Dk^2t}
\]

while expanding the radial part
\begin{align*}
  \pderiv[2]{\phi}{r} + \frac{2}{r}\pderiv{\phi}{r} + k^2\phi = 0\\
  r\pderiv[2]{\phi}{r} + 2\pderiv{\phi}{r} + k^2r\phi = 0\\
  \pderiv[2]{}{r}(r\phi) + k^2r\phi = 0
\end{align*}
whose solution is
\[
  \phi(r) = A\frac{\sin(kr)}{r} + B\frac{\cos(kr)}{r}
\]
applying the simmetry condition in $r = 0$ we have $B = 0$, while the boundary condition in $r = R$ gives
\[
  \sin(kR) = 0 \quad\Rightarrow\quad k_n = n\frac{\pi}{R}
\]
therefore the solution of the homogeneous problem will be in the form
\[
  w(r,t,s) = \sum_{n=1}^{\infty} b_n \frac{\sin(k_nr)}{r}e^{-Dk_n^2t}
\]
applying the initial condition we can find the Fourier coefficients $b_n$
\[
  y\dot{F}r = \sum_{n=1}^{\infty} b_n \sin(k_nr)
\]
\[
  b_n = \frac{2y\dot{F}}{R}\underbrace{\int_0^R r\sin(k_nr)}_{R(-1)^{n+1} / k_n} \, dr = 2y\dot{F}\frac{(-1)^{n+1}}{k_n}
\]

the solution of the homogeneous problem is
\[
  w(r,t,s) = 2y\dot{F}\sum_{n=1}^{\infty} \frac{(-1)^{n+1}}{k_n}\frac{\sin(k_nr)}{r}e^{-Dk_n^2t} \quad [\text{atomi}/m^3s]
\]

Finally, integrating we obtain the solution of the non-homogeneous problem
\[
  u(r,t) = \int_0^t w(r,t-s,s) \, ds = 2y\dot{F}\sum_{n=1}^{\infty} \frac{(-1)^{n+1}}{k_n}\frac{\sin(k_nr)}{r}\int_0^t e^{-Dk_n^2(t-s)} \, ds
\]
\[
  u(r,t) = \frac{2y\dot{F}}{D}\sum_{n=1}^{\infty} \frac{(-1)^{n+1}}{k_n^3}\frac{\sin(k_nr)}{r}\left(1 - e^{-Dk_n^2t}\right)\quad [\text{atomi}/m^3]
\]

To find the total number of atoms retained in the grain we integrate over the volume
\[
  g(t) = \int_{V_g} u(r,t) dV_g = 4\pi\int_0^R r^2 u(r,t) \, dr = \frac{8\pi y\dot{F}}{D}\sum_{n=1}^{\infty} \frac{(-1)^{n+1}}{k_n^3}\left(1 - e^{-Dk_n^2t}\right)\int_0^R r\sin(k_nr) \, dr
\]

\[
  g(t) = \frac{8y\dot{F}R^5}{\pi^3D}\sum_{n=1}^\infty \frac{1 - e^{-D\left( n\pi/R \right)^2t}} {n^4} \approx \frac{8y\dot{F}R^5}{\pi^3D}\left(1 - e^{-D\left( \pi/R \right)^2t} \right)
\]

\subsection{Release in the plenum}
The number of atoms released in the plenum by a single grain is given by
\[
  R_g(t) = p(t) - g(t) = y\dot{F}\left[ t - \frac{8yR^5}{\pi^3D}(1 - e^{-D\left( \pi/R \right)^2t})\right]
\]
the number of atoms released in the plenum by all the fuel
\[
  R_f(t) = R_g(t)N_g = R_g(t)\frac{V_f}{V_g}
\]

\end{document}
